\originalchapter{课程作业与中文题解}
本章按18.703 Spring 2013官方11份作业单的原编号收录完整公开题面. 官网未公开官方答案, 本章解答为AI独立撰写并数学交叉审读. 有书目出处但题面完整重印在作业单中的题仍保留出处; 只给外书题号的条目不补造题面. 字母 $C_m$ 表示阶 $m$ 循环群, $D_n$ 阶为 $2n$.

11份作业共126个原题号, 其中45题有完整公开题面, 69个明确小问; 其余81题只给书目定位. 模拟卷另有24题、47小问, 合计69个公开大题号、116小问 (去重前). 作业5第12题与第一次模拟期中第5题相同, 保留双映射; 作业5第9、11题是同一缺面定位, 故81个缺面题号只有80个不同定位. 其余近似主题保留各自原号. 计数把未分问的题视为一个小问, 不把定义中的列举条件另算小问.
\begin{longtable}{@{}p{.12\textwidth}p{.34\textwidth}p{.46\textwidth}@{}}
\toprule 作业&完整题面原题号&仅书目定位, 本册准确登记缺面\\\midrule\endhead
1&无&1--11\\
2&1,3,6&2,4,5,7--14\\
3&1,5,6&2--4,7\\
4&2,5&1,3,4,6--11\\
5&3,10,12--14&1,2,4--9,11\\
6&1,6,8--13&2--5,7\\
7&9&1--8,10--13\\
8&3,5,10&1,2,4,6--9,11--13\\
9&1--12&无\\
10&3,7&1,2,4--6,8\\
11&5--10&1--4\\
\bottomrule
\end{longtable}
逐题书名、章节、节和题号见随包 assessment-map.json; 对应官方原作业单链接可点击. 不拿当今 Judson 的同号题猜2013版题面, 不复制未核许可的 Herstein 外书题面. 作业1全部仅书目定位, 所以没有可诚实重构的题解. 下列章节保留原文件所列截止日期, 不改成新学期日历.
\originalsection{作业2: 2月28日}
\begin{exercise}原第1题. 令 $Z(G)=\{g:hg=gh\text{ 对所有 }h\in G\}$. (i) 证明 $Z(G)=\bigcap_{g\in G}C_G(g)$. (ii) 证明中心为子群.\end{exercise}
\begin{solution}
中心定义恰要求同时属于每个中心化子, 所以等式成立. 中心化子含单位, 且若 $x,y$ 与 $g$ 交换, 则 $xy^{-1}$ 也与 $g$ 交换, 故均为子群, 它们的交也是子群. 亦可直接用同一判别式验证中心.
\end{solution}
\begin{exercise}原第3题, 勘正条件. 非平凡群 $G$ 没有非平凡真子群, 证明它为素数阶循环群. 原题“没有真子群”须作此解释, 否则遗漏平凡群边界.\end{exercise}
\begin{solution}
任取 $a\ne e$, $\langle a\rangle$ 非平凡, 必为整个 $G$. 若无限阶, $\langle a^2\rangle$ 为非平凡真子群; 若有限合数阶 $n$, 取 $1<d<n$, $d\mid n$, 则 $\langle a^d\rangle$ 同样违背条件. 故阶为素数.
\end{solution}
\begin{exercise}原第6题. 找出正方形对称群 $D_4$ 的全部子群.\end{exercise}
\begin{solution}
写 $D_4=\langle r,s\mid r^4=s^2=e,srs=r^{-1}\rangle$. 子群为 $\{e\}$, $\langle r^2\rangle$, 四个 $\langle r^ks\rangle$ ($0\le k<4$), $\langle r\rangle$, $\langle r^2,s\rangle$, $\langle r^2,rs\rangle$, 全群, 共10个. 任意子群与旋转子群交为阶1、2或4的子群. 若还含翻折 $r^ks$, 任意另一个翻折与它的乘积是旋转, 因而子群由该交和这一翻折生成. 对三种交分别枚举正是上述列表, 无遗漏.
\end{solution}
\originalsection{作业3: 3月7日}
\begin{exercise}原第1题, Herstein第3章\S1第1、5题 (题面在作业单完整给出). 求下列乘积并求其阶, 从右向左复合:
\begin{align*}
\text{(i)}\quad&\begin{pmatrix}1&2&3&4&5&6\\6&4&5&2&1&3\end{pmatrix}
\begin{pmatrix}1&2&3&4&5&6\\2&3&4&5&6&1\end{pmatrix};\\
\text{(ii)}\quad&\begin{pmatrix}1&2&3&4&5\\2&1&3&4&5\end{pmatrix}
\begin{pmatrix}1&2&3&4&5\\3&2&1&4&5\end{pmatrix};\\
\text{(iii)}\quad&\alpha^{-1}\begin{pmatrix}1&2&3&4&5\\2&1&3&4&5\end{pmatrix}\alpha,
\quad\alpha=\begin{pmatrix}1&2&3&4&5\\4&1&3&2&5\end{pmatrix}.
\end{align*}
\end{exercise}
\begin{solution}
(i) 逐点先右后左得到 $1\mapsto4\mapsto1$, $2\mapsto5\mapsto3\mapsto2$, 6固定, 为 $(14)(253)$, 阶6. (ii) 为 $(132)$, 阶3. (iii) 换位共轭为 $(\alpha^{-1}(1)\ \alpha^{-1}(2))=(24)$, 阶2.
\end{solution}
\begin{exercise}原第5题. 对 $\sigma=(1472)(365)$, $\tau=(123)(475)$, 求 $\tau\sigma\tau^{-1}$ 及 $\sigma,\tau$ 的阶.\end{exercise}
\begin{solution}
将每个循环位置用 $\tau$ 的像替换, 得 $(2753)(164)$. 两原置换各不交循环长度为4与3, 3与3, 所以阶分别12与3.
\end{solution}
\begin{exercise}原第6题. 找 $\tau\in S_7$ 使 $\tau(125)(3674)\tau^{-1}=(314)(2765)$.\end{exercise}
\begin{solution}
按两循环的列出顺序要求 $1,2,5\mapsto3,1,4$, $3,6,7,4\mapsto2,7,6,5$. 得 $\tau=(132)(45)(67)$, 逐点代入共轭公式可核.
\end{solution}
\originalsection{作业4: 3月14日}
\begin{exercise}原第2题. 证明 $D_4$ 的旋转子群 $H=\{I,R,R^2,R^3\}$ 正规.\end{exercise}
\begin{solution}
$H$ 指数2, 由第2章的指数2判别正规. 也可对生成元检查: 旋转与旋转交换, 翻折将 $R^k$ 共轭为 $R^{-k}$, 均保持 $H$.
\end{solution}
\begin{exercise}原第5题. $G=S_3,H=\{e,(12)\}$. (i) 列所有左陪集. (ii) 列所有右陪集. (iii) 每个左陪集都是某个右陪集吗?\end{exercise}
\begin{solution}
左陪集为 $H,\{(13),(123)\},\{(23),(132)\}$; 右陪集为 $H,\{(13),(132)\},\{(23),(123)\}$. 逐个右乘或左乘 $(12)$ 可核. 含 $(13)$ 的左右两陪集不同, 且每个方向只有一个陪集包含该元素, 所以(iii)否.
\end{solution}
\originalsection{作业5: 3月21日}
\begin{exercise}原第3题, Herstein第2章\S6第3、4题. 设 $N\trianglelefteq G,\bar M\le G/N$, 令 $M=\{a:aN\in\bar M\}$. (i) 证明 $M\le G,N\subseteq M$. (ii) 若 $\bar M$ 正规则 $M$ 正规.\end{exercise}
\begin{solution}
$M$ 是商投影的原像, 非空且若 $a,b\in M$, $(ab^{-1})N=(aN)(bN)^{-1}\in\bar M$, 故成子群. $N$ 的像是单位, 所以包含 $N$. 若 $\bar M$ 正规, $gmg^{-1}$ 的像为 $(gN)(mN)(gN)^{-1}\in\bar M$, 得 $M$ 正规. 原卷把 $M$ 和 $\bar M$ 写成同一字母, 本册区分以免对象混淆.
\end{solution}
\begin{exercise}原第10题, Herstein第2章\S7第6题. 设 $N\trianglelefteq G$. 若 $a$ 阶为有限 $d$, 证明 $aN$ 的阶 $m$ 有限且整除 $d$.\end{exercise}
\begin{solution}
$(aN)^d=N$, 因而有限阶存在. 写 $d=qm+r$, $0\le r<m$, 得 $(aN)^r=N$, 最小性迫使 $r=0$.
\end{solution}
\begin{exercise}原第12题. 两正规子群交平凡, 证明它们逐元素交换.\end{exercise}
\begin{solution}
此题与第一次模拟期中第5题完全相同, 完整交换子论证统一见习题\ref{ex:commute}. 本题原编号保留, 不另计一份不同题面.
\end{solution}
\begin{exercise}原第13题. 证明 $G\cong H'\times K'$ 当且仅当 $G$ 有正规子群 $H\cong H',K\cong K'$, $H\cap K=\{e\}$ 且 $G=H\vee K$ (由两子群生成).\end{exercise}
\begin{solution}
上一题给出 $h,k$ 交换, 所以 $HK$ 是包含两者的子群, 就是生成子群 $H\vee K=G$. 映射 $H\times K\to G$, $(h,k)\mapsto hk$, 是满射同态, 核由 $h=k^{-1}\in H\cap K$ 知平凡. 反向在直积中取两个坐标子群, 它们正规、交平凡且生成全群, 沿同构搬回 $G$ 即得.
\end{solution}
\begin{exercise}原第14题 (挑战). 给有限集合上一种运算, 满足除结合律外的全部群公理.\end{exercise}
\begin{solution}
取五元素 $e,a,b,c,d$, 运算表为
\begin{center}\begin{tabular}{c|ccccc}
$\cdot$&$e$&$a$&$b$&$c$&$d$\\\hline
$e$&$e$&$a$&$b$&$c$&$d$\\
$a$&$a$&$e$&$c$&$d$&$b$\\
$b$&$b$&$d$&$e$&$a$&$c$\\
$c$&$c$&$b$&$d$&$e$&$a$\\
$d$&$d$&$c$&$a$&$b$&$e$
\end{tabular}\end{center}
表定义封闭运算, $e$ 为双侧单位元, 每个元素平方为 $e$ 所以有双侧逆元. 但 $(aa)b=b$, $a(ab)=ac=d$, 故不结合.
\end{solution}
\originalsection{作业6: 4月3日}
\begin{exercise}原第1题. 给二面体群 $D_n$ 的一个生成元关系表示.\end{exercise}
\begin{solution}
对 $n\ge3$, 表示为 $\langle r,s\mid r^n=s^2=e,srs=r^{-1}\rangle$. 关系将任何词整理为 $r^i$ 或 $r^is$, 至多 $2n$ 种, 正 $n$ 边形的具体旋转翻折满足这些关系且有 $2n$ 个不同变换, 所以所表示的群恰为 $D_n$.
\end{solution}
\begin{exercise}\label{ex:3generate}原第6题, Herstein第3章\S3第6题. 证明3循环生成 $A_n$.\end{exercise}
\begin{solution}
偶置换可按相邻两换位配对. 相同两换位抵消; 共一个字母的两换位之积是3循环; 不交情形有 $(ab)(cd)=(acb)(acd)$, 从右向左作用可逐点核实. 所以每对均可化为3循环的积, 而每个3循环本来是偶置换.
\end{solution}
\begin{exercise}原第8题. 证明 Sylow $p$ 子群为特征子群当且仅当它唯一.\end{exercise}
\begin{solution}
特征子群指被每个自同构保持, 特别被内自同构保持而正规, Sylow 共轭性给唯一. 反之自同构保持子群阶, 必置换所有 Sylow $p$ 子群; 唯一时就必须保持它.
\end{solution}
\begin{exercise}原第9题. 交换群 $G$ 阶为 $p^lq^m$, $p\ne q$, $l,m>0$. 证明 $G\cong P\times Q$, $|P|=p^l,|Q|=q^m$.\end{exercise}
\begin{solution}
取两 Sylow 子群, 交换性使它们正规, 阶互素使交平凡. 乘积子群阶 $|P||Q|=|G|$, 故生成全群. 作业5第13题的同构 $(a,b)\mapsto ab$ 给直积.
\end{solution}
\begin{exercise}原第10题. 证明阶168的单群同构于 $S_8$ 的子群.\end{exercise}
\begin{solution}
$n_7\mid24$, $n_7\equiv1\pmod7$, 只能1或8. 单性排除1. 共轭作用在八个 Sylow 7子群上是非平凡同态 $G\to S_8$, 核正规且不是全群, 故平凡, 得嵌入.
\end{solution}
\begin{exercise}\label{ex:an}原第12题 (挑战), Herstein第3章\S3第9、10题. (i) $n\ge5$, 若 $\{e\}\ne N\trianglelefteq A_n$, 证明 $N$ 含3循环. (ii) 在 $n\ge3$ 范围, 证明 $A_n$ 单当且仅当 $n\ne4$. 原题没有注明(ii)的低阶范围, $A_1,A_2$ 平凡不算单群.\end{exercise}
\begin{solution}
取 $\sigma\in N\setminus\{e\}$ 使移动点数 $m$ 最小. 若 $m\ge7$, 取 $a$ 使 $\sigma(a)=b\ne a$, 选不同的 $c,d$ 避开 $a,b$, 令 $\tau=(acd)$. $\sigma\tau\sigma^{-1}$ 的支撑含 $b$, $\tau$ 的支撑不含 $b$, 故不相等. 交换子 $\sigma\tau\sigma^{-1}\tau^{-1}\in N$ 非平凡, 最多移动两个三点支撑的并, 即6点, 矛盾. 所以 $3\le m\le6$.

$m=3$ 即3循环. $m=4$ 必为 $(ab)(cd)$; 选支撑外一点 $e_0$, 用 $\tau=(cde_0)$, 共轭为 $\tau^{-1}$, 交换子为 $\tau^{-2}=\tau$, 得3循环. $m=5$ 必为 $(abcde_0)$, 用 $\tau=(abc)$, 交换子 $(bcd)(acb)=(adb)$ 是3循环. $m=6$ 的偶置换类型仅 $(abcd)(e_0f)$ 或 $(abc)(de_0f)$. 前者同样取 $(abc)$ 得 $(adb)$; 后者取 $\tau=(abd)$, 共轭为 $(bce_0)$, 二者支撑不同, 交换子非平凡而至多移动5点, 与最小性矛盾. 这穷尽了无固定点的偶循环类型, 证明(i).

任意两个3循环在 $A_n$ 中共轭: 先在 $S_n$ 中逐位置匹配三点; 若匹配置换为奇, 再在目标循环补集两个点交换, 因 $n\ge5$ 存在这两点, 改为偶且保持匹配. 正规性使 $N$ 含全部3循环, 由习题\ref{ex:3generate}生成全 $A_n$, 得单性. $A_3=C_3$ 单, $A_4$ 的双换位子群 $V$ 非平凡真且正规, 故不单. 这给准确低阶边界.
\end{solution}
\begin{exercise}原第11题与第13题 (后者挑战). 分别证明阶在1至59、61至167的单群必为素数阶.\end{exercise}
\begin{solution}
先给统一约束. 交换单群必为素数阶循环群, 因每个子群正规且任意非单位元生成全群. 非交换单群 $G$ 对任何真子群 $H$ 的左陪集作用非平凡, 核正规故忠实. 符号同态必平凡, 否则 $G$ 嵌入阶2群; 因而 $G\hookrightarrow A_m$, $m=[G:H]$, 必须 $|G|\mid m!/2$. 同样, 若 Sylow 数 $n_p>1$, 共轭作用给 $G\hookrightarrow A_{n_p}$. 合数素数幂阶由中心非平凡排除: 单性使中心为全群, 会迫使交换, 只能素数阶.

下表对每个列出合数 $n$, 取 Sylow $p$ 子群 $H$, $m=n/p^{v_p(n)}$, 已直接核 $n\nmid m!/2$, 故排除. $v_p(n)$ 是整数分解中 $p$ 的指数. 表格是有限算术, 作用约束才是数学理由.
\begin{longtable}{@{}rrr@{\qquad}rrr@{\qquad}rrr@{}}
\toprule $n$&$p$&$m$&$n$&$p$&$m$&$n$&$p$&$m$\\\midrule\endhead
6&2&3&10&5&2&12&2&3\\
14&7&2&15&5&3&18&3&2\\
20&5&4&21&7&3&22&11&2\\
24&2&3&26&13&2&28&7&4\\
33&11&3&34&17&2&35&7&5\\
36&3&4&38&19&2&39&13&3\\
40&2&5&42&7&6&44&11&4\\
45&3&5&46&23&2&48&2&3\\
50&5&2&51&17&3&52&13&4\\
54&3&2&55&11&5&57&19&3\\
58&29&2&&&&&&\\
\bottomrule
\end{longtable}
1至59中余下合数仅30和56. 阶30单性迫使 $n_5=6,n_3=10$, 产生 $6\cdot4+10\cdot2=44$ 个不同非单位元, 超过29. 阶56单性迫使 $n_7=8$, 产生48个阶7元素; 剩余恰8元素, 任何阶8 Sylow 2子群都须等于此剩余集合, 因而唯一正规, 矛盾. 第11题得证.

61至167的陪集排除表如下:
\begin{longtable}{@{}rrr@{\qquad}rrr@{\qquad}rrr@{}}
\toprule $n$&$p$&$m$&$n$&$p$&$m$&$n$&$p$&$m$\\\midrule\endhead
62&31&2&65&13&5&66&11&6\\
68&17&4&69&23&3&74&37&2\\
75&5&3&76&19&4&77&11&7\\
78&13&6&80&2&5&82&41&2\\
85&17&5&86&43&2&87&29&3\\
88&11&8&91&13&7&92&23&4\\
93&31&3&94&47&2&95&19&5\\
96&2&3&98&7&2&99&11&9\\
100&5&4&102&17&6&104&13&8\\
106&53&2&108&3&4&110&11&10\\
111&37&3&112&2&7&114&19&6\\
115&23&5&116&29&4&117&13&9\\
118&59&2&119&17&7&122&61&2\\
123&41&3&124&31&4&129&43&3\\
130&13&10&133&19&7&134&67&2\\
135&3&5&136&17&8&138&23&6\\
141&47&3&142&71&2&143&13&11\\
145&29&5&146&73&2&147&7&3\\
148&37&4&150&5&6&152&19&8\\
153&17&9&155&31&5&156&13&12\\
158&79&2&159&53&3&160&2&5\\
161&23&7&162&3&2&164&41&4\\
166&83&2&&&&&&\\
\bottomrule
\end{longtable}
阶63,70,84,126,140,154,165分别取素数7,5,7,7,5,11,11, Sylow 整除和同余条件只允许 $n_p=1$, 所以不单. 阶72时 $n_3=1$ 或4; 后者迫使嵌入 $A_4$, 阶不容. 余下仅90,105,120,132,144, 逐一处理.

阶90: 单性给 $n_3=10,n_5=6$. 阶9 Sylow 子群交换. 若不同两者 $Q,Q'$ 交非平凡, 取交中的阶3元 $x$, 则两者都在 $C_G(x)$ 中. 此中心化子阶为9,18,45或90. 前三种阶的 Sylow 3子群都唯一 (个数整除1,2,5并模3余1); 最后一种使 $x$ 中央, 与非交换单性冲突. 所以两者交平凡. 产生 $10\cdot8+6\cdot4=104$ 个不同非单位元, 超过89.

阶105: 单性给 $n_7=15,n_5=21$, 两类阶素数元素共 $15\cdot6+21\cdot4=174$, 超过104. 阶120的完整排除见第二次模拟期中习题\ref{ex:120}; 其中所需 $A_6$ 单性已在本章前题证明. 阶132: $n_{11}=12$, $n_3=1,4$ 或22, 前两值被正规性或嵌入 $A_4$ 排除. 后值产生 $12\cdot10+22\cdot2=164$ 个不同非单位元, 超过131.

阶144: $n_3=1,4$ 或16, 单性和 $A_4$ 阶限制迫使16. 若不同阶9子群交非平凡, 对交中的阶3元 $x$, 中心化子阶为9,18,36,72或144. 前两阶只有一个 Sylow 3子群; 36,72的中心化子指数4,2, 陪集作用的忠实性不容阶144的单群; 144使 $x$ 中央. 所以16个 Sylow 3子群只交于单位, 给128个非单位奇数阶元素. 剩余恰16元素, 每个阶16 Sylow 2子群都等于剩余集合, 成为唯一正规子群, 矛盾. 两表、素数幂和这些余项覆盖区间全部合数, 第13题得证.
\end{solution}
\originalsection{作业7: 4月11日}
\begin{exercise}原第9题, Herstein第4章\S1第26题. 复四元数 $H(\C)$ 的元素为 $\alpha_0+\alpha_1i+\alpha_2j+\alpha_3k$, 系数在 $\C$, 四元数基元满足 $i^2=j^2=k^2=ijk=-1$. 证明它不是除环.\end{exercise}
\begin{solution}
用 $\iota\in\C$ 表示标量虚数, 避免与基元 $i$ 混淆. 两个非零元素 $1+\iota i$ 与 $1-\iota i$ 的乘积是 $1-(\iota i)^2=1-1=0$, 因标量与基元交换. 有非零零因子, 不可能每个非零元可逆: 若第一因子可逆, 左乘其逆会迫使第二因子为0.
\end{solution}
\originalsection{作业8: 4月18日}
\begin{exercise}原第3题. 对理想 $I,J$ 证明 (i) $I\cap J$ 为理想; (ii) $I+J=\{i+j\}$ 为理想; (iii) 有限和 $\sum i_kj_k$ 组成的 $IJ$ 为理想.\end{exercise}
\begin{solution}
交中元素的差以及左右乘法仍在两理想中. 和中 $(i+j)-(i'+j')=(i-i')+(j-j')$ 留在和中, 左右乘 $r$ 分别用 $ri,rj$ 或 $ir,jr$ 留在原理想. 积中和与负元逐项构造, 左乘将每项变成 $(ri_k)j_k$, 右乘变成 $i_k(j_kr)$, 仍为允许有限和. 这些证明对非交换环的双边理想也成立.
\end{solution}
\begin{exercise}原第5题. 固定素数 $p$, 令 $R$ 为既约分母不被 $p$ 整除的有理数. (i) 证明 $R\le\Q$ 为子环. (ii) 分子被 $p$ 整除的元素组成理想 $I$, 证明 $R/I\cong\F_p$.\end{exercise}
\begin{solution}
加乘运算后的分母可取原两分母的积, 不被 $p$ 整除, 约分也不引入 $p$, 且0,1及负元均在集合中. 映射 $a/b\mapsto\bar a\bar b^{-1}$ 良定义: 相等分数有 $ad=bc$, 在 $\F_p$ 中消去可逆分母. 映射保加乘及1, 整数像使它满射, 核恰是分子被 $p$ 整除的集合, 因而为理想, 第一同构定理给结论.
\end{solution}
\begin{exercise}原第10题, Herstein第4章\S3第22、23题. 互素正整数 $m,n$, 令 $I_m=m\Z,I_n=n\Z$. (i) 求交. (ii) 建立 $\Z/(mn)\to\Z/(m)\times\Z/(n)$ 的单射环同态. (iii) 证明它为同构.\end{exercise}
\begin{solution}
(i) 同时整除条件和互素性给交 $mn\Z$. (ii) 映射 $\bar a\mapsto(a\bmod m,a\bmod n)$ 保加乘, 良定义, 核由(i)为0. (iii) 取 $um+vn=1$, 对给定余数 $A,B$, $Avn+Bum$ 同时满足模 $m$ 为 $A$, 模 $n$ 为 $B$, 得满射. 若原模数为负, 理想不变, 改取绝对值; 模1允许零商环例外.
\end{solution}
原作业末尾还要求证明环直积有一般余积泛性质, 这个要求本身错误. 第5章练习以两恒等映射给出反例; 本册登记此附注并纠正, 不为假命题伪造证明.
\originalsection{作业9: 4月25日}
\begin{exercise}原第1题, Herstein第4章\S4第2题. 对高斯整数环, 分量均被3整除的集合 $M$ 为理想, 证明商为9元域.\end{exercise}
\begin{solution}
$M=(3)$, 所以是理想. 商同构于 $\F_3[x]/(x^2+1)$, 对应 $x\mapsto i$. 每个剩余类恰由两个模3分量确定, 共9个. $x^2+1$ 在 $\F_3$ 无根 (平方仅0,1), 故不可约, 商为域.
\end{solution}
\begin{exercise}原第2题, Herstein第4章\S4第3题. (i) 证明 $R=\Z[\sqrt2]$ 为复数子环. (ii) 分量均被5整除的 $M$ 为理想, 且商为25元域.\end{exercise}
\begin{solution}
乘积 $(a+b\sqrt2)(c+d\sqrt2)=(ac+2bd)+(ad+bc)\sqrt2$ 留在集合中, 加负元及1也成立. $M=(5)$, 商同构于 $\F_5[x]/(x^2-2)$, 共25个元素. 模5平方为0,1,4, 不含2, 二次式不可约, 故商域成立. 也可由共轭乘积 $a^2-2b^2$ 给非零类求逆.
\end{solution}
\begin{exercise}原第3题. 构造49元域.\end{exercise}
\begin{solution}
取 $\F_7[x]/(x^2+1)$. 模7平方为0,1,2,4, 不含 $-1=6$, 所以二次式不可约; 商域基为 $1,\bar x$, 有 $7^2=49$ 个元素.
\end{solution}
\begin{exercise}原第4题. 真理想 $I$ 外的每个元素均为单位, 证明 $I$ 为唯一极大理想.\end{exercise}
\begin{solution}
任一真理想不能含单位, 否则吸收其逆元得到1. 因此任何真理想均包含于 $I$, 其中 $I$ 本身为真理想, 所以既极大又唯一.
\end{solution}
\begin{exercise}原第5题. 环同态 $\varphi:R\to S$ 保1. (i) 素理想 $J\subset S$ 的原像为素理想. (ii) 给 $J$ 极大而原像非极大的例子.\end{exercise}
\begin{solution}
核验差和标量吸收即得原像为理想; 保1使1不在原像, 所以为真理想. 若 $ab$ 在原像, $\varphi(a)\varphi(b)\in J$, 素性给一因子在 $J$, 从而一因子在原像. 对(ii), 用 $\Z\hookrightarrow\Q$ 和 $J=(0)$, 域中的零理想极大, 原像为 $\Z$ 的零理想, 却严格位于 $(2)$ 下, 非极大.
\end{solution}
\begin{exercise}原第6题. 整环中证明 (i) $a\mid b\iff(b)\subseteq(a)$; (ii) $a,b$ 相伴当且仅当主理想相同; (iii) $a$ 为单位当且仅当 $(a)=R$.\end{exercise}
\begin{solution}
(i) $b=ac$ 使所有 $rb$ 是 $a$ 的倍数, 反向包含使 $b$ 本身为倍数. (ii) 相伴直接给相同理想; 反向写 $a=ub,b=va$, 若非零消去得 $uv=1$, 因交换环 $u$ 单位. 若一方为0, 理想相同迫使另一方也0, 二者相伴. (iii) $(a)=R$ 等价于 $1=ra$, 正是可逆性.
\end{solution}
\begin{exercise}原第7题. 证明整环的素元不可约.\end{exercise}
\begin{solution}
若素元 $p=ab$, 素性使 $p\mid a$ 或 $p\mid b$. 如 $a=pc$, 消去非零 $p$ 得 $1=cb$, 所以 $b$ 单位. 另一情况对称.
\end{solution}
\begin{exercise}原第8题. 在整环中, 同一对元素的两个 gcd 必相伴.\end{exercise}
\begin{solution}
每个 gcd 既是公因子又被所有公因子整除, 故二者互相整除. 第6题(ii)给相伴, 包括零情形.
\end{solution}
\begin{exercise}原第9题. 对 $R=\Z[\sqrt{-5}]$: (i) 证明 $2,3,1\pm\sqrt{-5}$ 不可约. (ii) 证明每个非零非单位有不可约分解. (iii) 证明非UFD.\end{exercise}
\begin{solution}
范数 $a^2+5b^2$ 乘法性成立, 单位恰 $\pm1$, 且范数2,3均无解. 四元素范数为4,9,6,6, 非平凡因子分解会要求范数2或3, 所以不可约. 任意可约非单位的两个非单位因子的范数均严格小于原范数, 按正整数范数归纳给有限分解. 最后 $6=2\cdot3=(1+\sqrt{-5})(1-\sqrt{-5})$, 不可约因子范数不同故无法相伴, 唯一性失败.
\end{solution}
\begin{exercise}原第10题. UFD中 (i) 证明 lcm 存在, 每个公倍数都被它整除, 且仅差相伴; (ii) 证明 gcd 与 lcm 的乘积与 $ab$ 相伴.\end{exercise}
\begin{solution}
对非零 $a,b$ 按全部不可约相伴类写指数 $r_i,s_i$, lcm 取最大指数, gcd 取最小指数. 整除等价于逐指数比较, 给存在、定义条件及唯一性. 因 $\min(r_i,s_i)+\max(r_i,s_i)=r_i+s_i$, 乘积与 $ab$ 相伴. 若一个为0, lcm为0、gcd取另一元; 两者都0时均取0, 条件仍成立.
\end{solution}
\begin{exercise}原第11题 (挑战). 在有单位元的交换半群中定义唯一分解.\end{exercise}
\begin{solution}
原题的“半群”实际含单位元, 本册称交换幺半群. 单位 $u$ 指存在 $v$ 使 $uv=1$. 相伴指相差一个单位因子. 非单位 $a$ 为原子指 $a=bc$ 迫使某因子为单位. 唯一分解指每个非单位为有限原子乘积, 且任意两个原子分解的原子数相同, 经重排后逐项相伴. 单位以空积和一个单位系数处理. 一般不消去半群可能有异常, 不能无条件套用整环证明; 下一题在整数向量子幺半群中可消去.
\end{solution}
\begin{exercise}原第12题 (挑战). 给 $v_1,\ldots,v_n\in\Z^2$, 令 $S$ 是其非负整数线性组合, 运算为加法. 确定哪些 $S$ 有唯一分解.\end{exercise}
\begin{solution}
必须先处理单位: $U=S\cap(-S)$ 是单位群, $s,t$ 相伴当且仅当 $s-t\in U$. 在约化幺半群 $\bar S=S/U$ 中, 取全部不同原子类 $\bar a_1,\ldots,\bar a_r$. 精确判据为这些类在群 $\operatorname{gp}(S)/U$ 中无非零整数关系; 等价于映射 $\N^r\to\bar S$, $(n_i)\mapsto\sum n_i\bar a_i$, 为双射. $\operatorname{gp}(S)$ 表示生成的整数加法群.

说明原子确实有限且生成. 令 $C$ 为原生成向量的非负实锥, $L=C\cap(-C)$ 为其直线部分. 有 $L=\operatorname{span}_{\R}U$. 一方面单位方向显然在 $L$; 另一方面, 位于 $L$ 的生成向量的负方向可由原生成向量作非负有理组合 (整数系数线性系统存在实非负解时可取有理解). 清分母使 $-kv\in S$, 加上 $(k-1)v$ 得 $-v\in S$, 该生成向量为单位. 直线部分由落在其中的生成向量张成: 若 $x,-x\in C$, 将两表达相加为0, 参与该零和的生成方向均有负向在锥中. 因此两边实张成相同.

在 $\R^2/L$ 中, 非单位生成方向组成有限尖锥, 可选有理线性函数 $\ell$, 对它们严格正且在 $L$ 为0. 一维时它们同向; 二维时尖锥的两个极端方向夹角小于半周, 取其内侧法向之间的有理方向即可使内积全正. 清分母使 $\ell(S)$ 为非负整数, 且非单位取正值. 若非单位分成两个非单位, 两者次数 $\ell$ 均下降, 按次数归纳得到有限原子分解. 原子在原生成向量表达中只能有一个非单位生成元, 否则可拆, 其余都是单位. 所以原子类来自有限原列表, 删去单位、重复伴随类及可分解者后得到全部原子类.

这些类生成 $\bar S$, 满射性已有. 两个非负指数组合相同恰是两个不同原子分解, 所以单射性恰为唯一性. 非零整数关系拆成正负两边便产生不同分解; 它不能只含同号项, 因非单位之和不可能是单位: 若 $a+b+c=0$, 则 $b+c$ 是 $a$ 的逆, 使 $a$ 已是单位.

向量和 $U$ 的基为整数, 在商实空间的线性关系可取有理系数, 清分母并再乘一个整数使落入 $U$ 的格, 故整数无关系等价于 $a_i$ 在 $\R^2/\operatorname{span}_{\R}U$ 中线性无关. 因而分类为: $\operatorname{rank}U=0$ 时, $S=\{0\}$ 或由一个非零向量或两个不共线向量自由生成; $\operatorname{rank}U=1$ 时, $S=U$ 或 $S=U+\N a$ 且 $a\notin\operatorname{span}_{\R}U$; $\operatorname{rank}U=2$ 时 $S=U$, 全是单位. 例如 $(2,0),(0,3)$ 生成的幺半群唯一分解, 无须是整个 $\Z^2$ 的格基; $(2,0),(3,0)$ 则有 $3(2,0)=2(3,0)$, 不唯一.
\end{solution}
\originalsection{作业10: 5月2日}
\begin{exercise}原第3题. 求 $\Z[i]$ 中 $11+7i$ 与 $8-i$ 的 gcd.\end{exercise}
\begin{solution}
作欧几里得除法 $11+7i=(1+i)(8-i)+2$, 再 $8-i=4\cdot2-i$. 余项 $-i$ 为单位, 故 gcd 与1相伴. 可复核的贝祖式为 $1=-4i(11+7i)+(-4+5i)(8-i)$.
\end{solution}
\begin{exercise}原第7题 (挑战). 给一个环和元素 $a$, 它存在不可约分解, 但从 $a$ 开始的因子分解算法有不停止的选择路径.\end{exercise}
\begin{solution}
取独立变量 $u,v,w$ 和 $A=\Q[u,v,w,w^{-1}]$, 定义子环
\[
R=\Q[u,v,w,uv/w,uv/w^2,uv/w^3,\ldots].
\]
它为整环. 其中单项式 $u^av^bw^c$ 的指数 $a,b\ge0$, 且 $c<0$ 时必须 $a,b\ge1$; 反向每个这样的单项式都由所列生成元得到, 这是检查因子属于 $R$ 的依据.

$A$ 是将 UFD $\Q[u,v,w]$ 中的 $w$ 变为单位得到的环, 仍唯一分解: 清掉 Laurent 分母后在多项式环中分解, 只把 $w$ 的幂并入单位. 其单位恰为 $\lambda w^j$, $\lambda\in\Q^\times,j\in\Z$, 因清分母后的可逆性只容 $w$ 因子. 所以 $u=fg$ 在 $A$ 中, 除交换次序外必为 $f=\lambda uw^j,g=\lambda^{-1}w^{-j}$. 两者若属于 $R$, 后者迫使 $j\le0$, 前者因无 $v$ 而迫使 $j\ge0$, 故 $j=0$, 某因子单位. 因而 $u$ 是 $R$ 原子, $v$ 同理, $a=uv$ 有原子分解.

另一方面令 $a_k=uv/w^k$, 则 $a_k=w a_{k+1}$ 对所有 $k\ge0$ 成立. 所有因子属于 $R$ 且非单位: $a_k$ 在 $A$ 已非单位, $w$ 的逆 $w^{-1}$ 不属于 $R$. 永远选 $a_{k+1}$ 继续拆便无穷, 每步都是非单位分解. 这不否认 $a_0$ 存在另一条有限分解, 也不宣称 $R$ 每个元素都有原子分解.
\end{solution}
\originalsection{作业11: 5月9日}
官网此文件是第11份作业, 原 PDF 标题误写“作业9”; 以下按官方目录顺序标11, 保留日期对应.
\begin{exercise}原第5题. $F[x]$ 自同构 $\varphi$ 固定每个 $a\in F$. (i) 证明 $f$ 不可约当且仅当 $\varphi(f)$ 不可约. (ii) 证明次数保持.\end{exercise}
\begin{solution}
(i) 自同构双向保持乘积和单位, 把非平凡分解双向搬运, 故保持不可约. (ii) 令 $h=\varphi(x),k=\varphi^{-1}(x)$, 它们都非恒定, 否则逆映射固定常数会使 $x$ 为常数. $k(h(x))=x$, 域上非恒定多项式复合次数为次数之积, 得 $\deg h=\deg k=1$. 因 $\varphi$ 固定系数, $\varphi(f)=f(h)$, 对非零 $f$ 次数保持; 零多项式双方同为0, 其次数按不定义或 $-\infty$ 的一致约定处理.
\end{solution}
\begin{exercise}原第6题. $b\in F^\times,c\in F$, 证明 $f(x)\mapsto f(bx+c)$ 为固定 $F$ 的自同构.\end{exercise}
\begin{solution}
代入保持和、积及常数, 逆代入为 $x\mapsto b^{-1}(x-c)$, 两次复合均为 $x$, 所以是自同构.
\end{solution}
\begin{exercise}原第7题. 证明所有固定 $F$ 的 $F[x]$ 自同构都是上述形式.\end{exercise}
\begin{solution}
由第5题(ii), $\varphi(x)=bx+c$ 且 $b\ne0$. 固定每个系数给 $\varphi(\sum a_ix^i)=\sum a_i(bx+c)^i$, 完成分类.
\end{solution}
\begin{exercise}原第8题. (i) 给 $\Q[x]$ 非恒等且平方为恒等的自同构. (ii) 对任意 $n>0$, 给 $\C[x]$ 阶为 $n$ 的自同构.\end{exercise}
\begin{solution}
(i) 取 $x\mapsto-x$. (ii) 取 $x\mapsto\zeta_n x$, $\zeta_n=\exp(2\pi i/n)$ 是本原 $n$ 次单位根, 固定复系数. 第 $k$ 次幂送 $x$ 至 $\zeta_n^kx$, 恰当 $n\mid k$ 才恒等, 所以阶为 $n$, 包括 $n=1$ 的恒等映射.
\end{solution}
\begin{exercise}原第9题. 对特征 $p>0$ 的域: (i) 证明 $(a+b)^p=a^p+b^p$. (ii) 证明 $q=p^n$ 次幂映射为环同态. (iii) 证明单射. (iv) 有限域时为自同构.\end{exercise}
\begin{solution}
特征 $p$ 必素, 二项式中间系数被 $p$ 整除, 给(i). 迭代 $n$ 次得 $q$ 次幂加法性, 乘法与1保持直接成立. 若 $a^q=b^q$, 则 $(a-b)^q=0$, 域无非零幂零元, 得 $a=b$. 有限集单射为满射, 故(iv).
\end{solution}
\begin{exercise}原第10题 (挑战). 给特征 $p$ 的域使 $a\mapsto a^p$ 不满射.\end{exercise}
\begin{solution}
取 $F=\F_p(t)$. 若 $t=(f/g)^p$, 清分母得 $f^p=tg^p$, 不可约因子 $t$ 的指数左边是 $p$ 倍数, 右边是 $1$ 加 $p$ 倍数, 不可能. 因而 $t$ 没有原像.
\end{solution}
